\section{Conclusiones}
A partir de las experiencias realizadas, se concluye que la variación de los parámetros de resistencia y capacitancia modifica significativamente la respuesta transitoria del circuito RLC serie. En particular, al aumentar la resistencia se observa un comportamiento más amortiguado, mientras que al disminuirla el sistema presenta oscilaciones antes de estabilizarse. Este comportamiento resulta consistente con el marco teórico.

Por otra parte, se verifica que, en un RLC serie, al disminuir la capacitancia también se reduce el tiempo necesario para que el circuito alcance el equilibrio. Esta observación concuerda con las predicciones teóricas, ya que la capacitancia influye en la dinámica del sistema a través de la frecuencia natural y, en consecuencia, en su tiempo característico de estabilización.

También, comprobamos la validez de la aproximación de que $\tau \sim \frac{1}{s_\text{dom}}$  en un RLC serie. En adición, comprobamos la validez de la aproximación de $V_\text{out} \sim V(5\tau)$ tanto para un RLC serie como para un RL. 



Otras conclusiones interesantes del informe...
\newpage